% ============================================
%  ВВЕДЕНИЕ
% Введение к ВКР, независимо от уровня подготовки,
% включает в себя следующие элементы:
% ˗ актуальность темы исследования;
% ˗ степень ее разработанности;
% ˗ решаемую проблему;
% ˗ цель и задачи;
% ˗ практическую значимость работы.
% ============================================

Автономные устройства с голосовым управлением~--- носимая электроника,
датчики промышленного мониторинга, беспроводные узлы умного дома~---
сталкиваются с общим архитектурным ограничением: конечным запасом энергии автономного источника.
Микрофон должен функционировать в режиме непрерывного мониторинга, иначе устройство
пропустит голосовую команду, однако постоянная работа аудиопайплайна на
микроконтроллере приводит к разряду источника питания в течение единиц
часов вместо месяцев и лет, ожидаемых от энергоэффективных
IoT-устройств. Указанное противоречие между непрерывной готовностью к
приёму команды и ограниченным энергетическим бюджетом определяет
актуальность исследований в области энергоэффективных систем
распознавания ключевых слов на периферии сети.

Альтернативой локальному распознаванию служат облачные голосовые
ассистенты (Amazon Alexa, Google Assistant, Apple Siri), передающие
аудиосигнал на удалённые серверы. Для батарейных устройств такой подход
неприемлем по совокупности факторов: непрерывная передача аудиоданных по
беспроводному каналу потребляет существенно больше энергии, чем локальная
обработка~\cite{src:9}; устройство становится зависимым от наличия и
качества сетевого соединения; передача голосовых данных создаёт риски для
конфиденциальности пользователя. Единственным практически реализуемым
вариантом для батарейных IoT-устройств остаётся выполнение распознавания
ключевых слов (\textit{Keyword Spotting}, KWS) непосредственно на
микроконтроллере~--- так называемый \textit{edge KWS}, распознавание на
периферии сети.

Задача распознавания ключевых слов на микроконтроллерах активно
исследуется международным научным сообществом, и к настоящему моменту в
литературе выделились три основных направления повышения
энергоэффективности. Первое направление~--- разработка компактных
нейросетевых архитектур, адаптированных к ресурсным ограничениям
встраиваемых платформ~\cite{src:1, src:12}. Второе~--- применение методов
квантизации весов и активаций, обеспечивающих переход от вычислений с
плавающей точкой к целочисленным операциям с использованием SIMD-расширений
современных микроконтроллеров~\cite{src:3, src:22}. Третье~--- построение
каскадных схем детекции, в которых вычислительно затратный нейросетевой
классификатор активируется не непрерывно, а лишь по событию,
зафиксированному более дешёвым предварительным детектором речевой
активности~\cite{src:6, src:7}. Вопросы построения детекторов голосовой
активности и компрессии моделей рассмотрены и в отечественной
литературе~\cite{src:23, src:24, src:26, src:27}. Однако перечисленные
направления в известной литературе изучаются преимущественно в изоляции
друг от друга: их совместное применение на микроконтроллере общего
назначения с количественной декомпозицией итогового энергобюджета по
этапам обработки не охарактеризовано. Восполнение указанного пробела и
составляет научно-техническую проблему настоящей работы.

Практический выбор, перед которым оказывается разработчик автономного
устройства с голосовым управлением, сводится к двум крайностям.
Микроконтроллер общего назначения в режиме непрерывного мониторинга
обеспечивает гибкость, доступность и развитую экосистему разработки, но
демонстрирует потребление, неприемлемое для батарейного питания.
Специализированный нейросетевой процессор (ASIC) демонстрирует
энергоэффективность, превосходящую микроконтроллер на один-два порядка,
но представляет собой закрытое решение с аппаратно фиксированной
архитектурой модели, ограниченно доступное и не допускающее независимой
модификации. Промежуточный вариант~--- каскадная архитектура с
квантизованной моделью на микроконтроллере общего назначения~--- известен
концептуально, однако количественно не охарактеризован: раздельный вклад
каскадного управления режимами питания и квантизации модели в итоговое
энергопотребление, как и вклад сервисных этапов цикла обработки
(пробуждение, инициализация периферии, захват аудио, индикация
результата) в общий бюджет энергии~--- не определены.

Целью настоящей работы является снижение энергопотребления устройств
интернета вещей с функцией голосового управления за счёт каскадного
управления режимами работы микроконтроллера и применения квантизованных
нейронных сетей для распознавания ключевых слов голосовых команд.
Достижение указанной цели предполагает решение следующих задач:

\begin{enumerate}
	\item проанализировать существующие методы и системы распознавания
	      ключевых слов голосовых команд, выявить основные источники
	      избыточного энергопотребления и сформулировать количественные
	      требования к энергоэффективному решению по совокупности метрик
	      (средняя мощность, энергия на детекцию, латентность, точность
	      классификации, размер модели);
	\item разработать каскадную архитектуру детекции из трёх этапов
	      возрастающей вычислительной сложности с согласованным
	      управлением режимами энергопотребления микроконтроллера,
	      допускающую перенос на различные микроконтроллерные платформы;
	\item исследовать влияние методов квантизации~--- посттренировочной
	      (PTQ) и квантизации с дообучением (QAT)~--- на точность,
	      размер и вычислительную сложность нейросетевой модели
	      классификатора ключевых слов;
	\item реализовать разработанную систему на платформе ESP32-S3 и
	      провести её экспериментальную апробацию на лабораторном стенде
	      с прецизионным измерением тока и напряжения;
	\item выполнить декомпозицию итогового энергобюджета цикла каскадной
	      обработки по этапам и дать количественную оценку вклада каждого
	      компонента архитектуры в общую экономию энергопотребления.
\end{enumerate}

\textit{Объектом} исследования является система распознавания ключевых
слов голосовых команд на микроконтроллере общего назначения;
\textit{предметом}~--- методы снижения её энергопотребления за счёт
каскадной архитектуры детекции и квантизации нейросетевой модели.

Научная новизна работы заключается в том, что впервые получено
количественное описание совместного применения каскадной архитектуры и
квантизации нейросетевой модели на микроконтроллере общего назначения с
раздельной оценкой вклада каждого метода в итоговое снижение
энергопотребления. В отличие от существующих работ, рассматривающих
указанные подходы изолированно, проведена покомпонентная декомпозиция
энергобюджета цикла обработки на платформе ESP32-S3, позволяющая
идентифицировать этапы с доминирующим вкладом в общий бюджет и
определить приоритетные направления дальнейшей оптимизации. В рамках
исследования впервые проведено сравнение методов посттренировочной
квантизации и квантизации с дообучением применительно к архитектуре
DS-CNN на платформе ESP32-S3 с одновременной оценкой влияния метода
квантизации на точность модели, её размер и латентность инференса в
условиях реального аппаратного ускорения.

На защиту выносятся следующие положения:

\begin{enumerate}
	\item Каскадная архитектура с тремя ступенями возрастающей
	      вычислительной сложности и согласованным управлением режимами
	      питания микроконтроллера обеспечивает кратное снижение средней
	      потребляемой мощности относительно режима непрерывного
	      мониторинга при сохранении точности детекции, что подтверждено
	      экспериментально на платформе ESP32-S3;
	\item Квантизация модели DS-CNN до целочисленного представления INT8
	      обеспечивает сжатие модели и сокращение времени инференса при
	      потере точности классификации, не превышающей статистической
	      погрешности оценки на тестовой выборке, и применима в обеих
	      формах (PTQ и QAT) к задаче распознавания ключевых слов на
	      ресурсоограниченных платформах;
	\item Покомпонентная декомпозиция энергобюджета каскадного цикла,
	      выполненная средствами прецизионных измерений тока и
	      напряжения, позволяет выделить этап с доминирующим вкладом в
	      общий бюджет энергии и определить приоритеты дальнейшей
	      оптимизации системы.
\end{enumerate}

Практическая значимость работы состоит в том, что предложенная
каскадная архитектура и методика прецизионных измерений
энергопотребления применимы при проектировании автономных IoT-устройств
с голосовым управлением в задачах промышленного мониторинга (голосовое
управление в условиях занятых рук оператора), носимой электроники
(медицинские и спортивные устройства) и беспроводных узлов умного дома.
Архитектура не привязана к конкретной аппаратной платформе и переносима
на любой микроконтроллер с поддержкой аппаратных режимов глубокого сна
и внешнего пробуждения~--- семейства STM32, nRF52 и аналогичные.
Результаты сравнительного исследования методов квантизации служат
ориентиром при выборе стратегии оптимизации нейросетевых моделей для
периферийных устройств, а методика двухдоменного измерительного стенда
с изоляцией измеряемой системы и системы измерений применима для
характеризации энергопотребления произвольных встраиваемых платформ.
